\documentclass[10pt,a4paper]{amsart}
\usepackage[margin=19mm]{geometry}
\usepackage{amsmath,amssymb,mathtools}
\usepackage[scr=rsfs,cal=euler]{mathalfa}
\usepackage{xfrac}
\usepackage[unicode,hidelinks,bookmarks=false]{hyperref}
\newcommand{\ie}{i.e.\ }
\newcommand{\cf}{cf.\ }
\newcommand{\resp}{respectively\ }
\newcommand{\Subsection}{\subsection}
\providecommand{\nobreakdash}{\nobreak\hbox{-}}
\setlength{\emergencystretch}{2em}
\multlinegap0pt
\usepackage{bm}
\usepackage{bbm}
\usepackage{dsfont}
\usepackage{xspace}
\usepackage{cancel}
\usepackage{mathtools}
\usepackage{amsthm}
\makeatletter
\@ifundefined{theorem}{\newtheorem{theorem}{Theorem}}{}
\@ifundefined{lemma}{\newtheorem{lemma}[theorem]{Lemma}}{}
\@ifundefined{proposition}{\newtheorem{proposition}[theorem]{Proposition}}{}
\makeatother
\newcommand{\RR}{\mathbb{R}}
\newcommand{\vare}{{\varepsilon}}
\newcommand{\vphi}{{\varphi}}
\renewcommand{\tilde}{\widetilde}
\title[Full range of infinite point blow-up exponents for the criti]{Full range of infinite point blow-up exponents for the critical generalized KdV equation}
\author{Nailya Manatova}
\date{Source: arXiv:2511.13538v1 (CC BY 4.0)}
\begin{document}
\maketitle

\noindent\emph{Source and licence.} Full range of infinite point blow-up exponents for the critical generalized KdV equation. Version arXiv:2511.13538v1 (CC BY 4.0), distributed under the Creative Commons Attribution 4.0 International licence, as recorded in the arXiv licence metadata for this exact version and shown on the abstract page https://arxiv.org/abs/2511.13538v1. Full rights notice: licenses/arxiv-2511.13538-CC-BY-4.0.txt.\par\smallskip

\noindent\textit{Editorial scope.} Counted unit: the Proposition whose source label is \texttt{Prop\_on\_blow\_up} in the article (source0.tex lines 2559--2580, source label \texttt{Prop\_on\_blow\_up}), delivered together with the article's own complete original proof of it (source0.tex lines 2604--2782, one \texttt{proof} environment of 179 lines). No mathematics is written, repaired or re-derived: statement and proof are the article's own text line for line. The proof is detached from the statement in the source: the article states the result at lines 2559--2580 and prints its proof at lines 2604--2782, and the proof environment's own header names the statement, so the association is the article's own. Required context retained uncounted: gKdV\_principal\_eq (lines 238--240); def\_of\_W\_and\_of\_r (lines 792--794); condition\_on\_rho (lines 797--799); BS1 (lines 802--808); lemma\_on\_decomposition\_around\_Q (lines 1298--1318); estimate\_on\_L2\_norm\_of\_vare (lines 1483--1491); estimate\_on\_L2\_norm\_of\_grad\_vare (lines 1488--1490); BS2 (lines 1719--1721); conseq\_of\_BS\_on\_g\_and\_hs (lines 1740--1742); Lyapounov\_control\_of\_H (lines 1805--1811); coercivity\_of\_mathcal\_F (lines 1813--1815); BS3 (lines 2553--2555). Every definition, display and label of those lines is retained unchanged. Omitted from this package and not redistributed: 1--237, 241--791, 795--796, 800--801, 809--1297, 1319--1482, 1491--1718, 1722--1739, 1743--1804, 1812--1812, 1816--2552, 2556--2558, 2581--2603, 2783--2857. Nothing inside a retained excerpt is omitted or altered: every retained body is delivered exactly as the article's lines are written, so no omission receipt of any kind accompanies this package. Citations: the retained excerpts carry no citation command at all, so no bibliography entry of the article is transcribed and no reference is redistributed. Reference adaptation: none was needed and none was made; every reference inside the delivered unit resolves inside it, so no unresolved reference or citation is produced. Printed slips kept literal: no printed word, symbol, hypothesis or number of the counted statement or of its proof was altered. Layout and class: the article's own class is \texttt{smfart} with packages including fontenc, babel, inputenc, xcolor, appendix, lipsum, amsmath, amssymb, amsfonts, smfthm, mathrsfs, systeme, stmaryrd, pgfplots; the wrapper is the lane's pinned \texttt{amsart} editorial wrapper at 10pt on A4, which restates the theorem-like environments the retained text needs and adds the article's own macro definitions that the retained text needs (\texttt{\textbackslash RR}, \texttt{\textbackslash tilde}, \texttt{\textbackslash vare}, \texttt{\textbackslash vphi}), with extra wrapper packages (bm, bbm, dsfont, xspace, cancel, mathtools). The delivered numbering is the unit's own. Assets: no retained span contains a figure environment, a graphics-inclusion command, a PDF-page inclusion, a table or a TikZ picture; no image file is copied into this package, the package declares no asset dependency and no image file is redistributed with it. Licence: version arXiv:2511.13538v1 of this article is distributed under the Creative Commons Attribution 4.0 International licence, as recorded in the arXiv licence metadata for this exact version and shown on the abstract page https://arxiv.org/abs/2511.13538v1. A full rights notice is delivered at licenses/arxiv-2511.13538-CC-BY-4.0.txt. Characters: every retained excerpt is pure ASCII, so the delivered engine, XeLaTeX with its default Latin Modern text font, reports no missing character.
\begin{equation}\label{gKdV_principal_eq}
    \partial_t U + \partial_x(\partial_{xx}U + U^5)=0, \qquad (t,x)\in [0,T) \times\RR.
\end{equation}

\begin{equation}\label{def_of_W_and_of_r}
    W(s,y) = Q_{b(s)}(y)+r(s)R(y), \qquad \text{with} \qquad r(s):=F(s,0)= \lambda^{\frac{1}{2}}(s)f(\tau(s),\sigma(s))
\end{equation}

\begin{equation}\label{condition_on_rho}
0<\rho < \frac{1}{4} \min{\{ 2\beta-1,\, \frac{1}{4}\, ,2-2\beta \}}.
\end{equation}% 

\begin{equation}\tag{BS1}\label{BS1}
    \begin{cases}
        \lvert \sigma(s) - \frac{1}{1-\beta} s^{1-\beta}\rvert \;\leq\; s^{1-\beta-\rho} \,, \\
        \lvert \lambda(s) - s^{-\beta} \rvert \;\leq\; s^{-\beta-\rho} \,,\\
        \lvert b(s) - \beta s^{-1} \rvert \;\leq\; s^{-1-\rho} \,,
    \end{cases}
\end{equation}

\begin{lemma}\label{lemma_on_decomposition_around_Q}
    Assume $v$ is close to the approximate blow-up profile for some $\alpha^{*}$ sufficiently small. There exist unique continuous functions $(\lambda,\sigma,b):I\mapsto (0,+\infty) \times \RR^2$ such that
    \begin{equation}\label{decomposition_of_v_in_lemma}
        \lambda^{\frac{1}{2}}(t)v(t, \lambda(t)y+\sigma(t)) = W(t,y)+ \vare(t,y)
    \end{equation}
    where $W$ is defined in \eqref{def_of_W_and_of_r} and  $\vare$ satisfies for all $t\in I$
    \begin{equation}\label{orthogonal_conditions_in_lemma}
        (\vare(t),y\Lambda Q ) = (\vare(t),\Lambda Q) = (\vare(t), Q) = 0.
    \end{equation}
    Moreover,
    \begin{equation}
        \bigg\lvert \frac{\lambda(t)}{\lambda_{\sharp}(t)} - 1\bigg\rvert + |b(t)| + \bigg\lvert \frac{\sigma(t)-\sigma_{\sharp}(t)}{\lambda_{\sharp}(t)} \bigg\rvert \lesssim \|z(t)\|_{L^2_{sol}}
    \end{equation}
    and
    \begin{equation}\label{estimates_on_norms_of_vares_in_z_in_decomp}
        \|\vare(t)\|_{L^2} \lesssim \|z(t)\|^{\frac 58}_{L^2}, \qquad \|\vare(t)\|_{L^2_{sol}} \lesssim \|z(t)\|_{L^2_{sol}},
    \end{equation}
    \begin{equation}\label{estimates_on_norms_of_partial_vares_in_z_in_decomp}
        \|\partial_y \vare(t)\|_{L^2} \lesssim \|z(t)\|_{L^2_{sol}} + \|\partial_y z(t)\|_{L^2},  \qquad \|\partial_y \vare(t)\|_{L^2_{sol}} \lesssim \|z(t)\|_{L^2_{sol}}+\|\partial_y z(t)\|_{L^2_{sol}}.
    \end{equation}
\end{lemma}  

\begin{lemma}[Mass and energy estimates]
    Under the bootstrap assumption \eqref{BS1}, for $s_0 \gg 1$, the following holds true on $s(I)$
    \begin{equation}\label{estimate_on_L2_norm_of_vare}
        \|\vare(s) \|^2_{L^2} \lesssim \Big\lvert\int U^2_0 - \int Q^2  \Big\rvert + s^{-2+\gamma}+ x^{-2\theta+1}_0,
    \end{equation}
    \begin{equation}\label{estimate_on_L2_norm_of_grad_vare}
        \lambda^{-2}\|\partial_y \vare\|^2_{L^2} \lesssim \lvert E(U_0)\rvert+s^{-3+4\beta}+x^{-2\theta-1}_0 + |g(s)|.
    \end{equation}
\end{lemma}

\begin{equation}\tag{BS2}\label{BS2}
        \mathcal{N}_B(s) \;\leq\; s^{-\frac{5}{4}}\,, \qquad \|\vare(s)\|_{H^1} \;\leq\; \alpha^*. 
\end{equation}

\begin{equation}\label{conseq_of_BS_on_g_and_hs}
    \lvert g(s) \rvert \lesssim \lvert g(s_0) \rvert  + s^{2\beta - 1-4\rho} \qquad \text{and}\qquad \lvert h_s(s) \rvert \lesssim s^{-\frac{\beta}{2}-1-4\rho},
\end{equation}

    \begin{equation}\label{Lyapounov_control_of_H}
    \begin{split}
        & \frac{d}{ds}\big[\mathcal{H}\big]+s^{j}\frac{\mu}{8}\int\vphi'_B(\vare_y^2+\vare^2) +2s^{j}\int_{y<-\frac B2}\psi'_B\,(\partial^2_y\vare)^2\\
        & +s^{j}\int_{y<-\frac B2} \psi'_B\,(\partial_y \vare)^2 +\frac{\lambda^k}{16}\int \vphi'_k\,\big(\vare^2 + (\partial_y\vare)^2 \big)
        \lesssim \lambda^ks^{ - \frac 52} + s^{j-4}.
    \end{split}
    \end{equation}

    \begin{equation}\label{coercivity_of_mathcal_F}
        \mathcal{F}\;\gtrsim\;\mathcal{N}^2_B.
    \end{equation}

\begin{equation}\tag{BS3}\label{BS3}
    \lvert h(s) \rvert \leq s^{-\frac{\beta}{2} - 2\rho}.
\end{equation}

\begin{proposition}\label{Prop_on_blow_up}
Let $x_0$ and $s_0>0$ be large enough and fix $\sigma_0$, $b_0$ such that
\begin{equation}\label{init_cond_sigma_b}
    \Big\lvert \sigma_0 - \frac{1}{1-\beta}s^{1-\beta}_0 \Big\rvert \leq s_0^{1-\beta-2\rho} \qquad \text{and} \qquad \Big\lvert b_0 - \beta s_0^{-1} \Big\rvert \leq s_0^{-1-2\rho},
\end{equation}
where $\rho $ satisfies \eqref{condition_on_rho}.
Let $\vare_0 \in H^1(\RR)$ be such that 
\begin{equation}\label{conditons_ortho_and_decroissance_on_vare_0}
s_0^{j} \, \|\varepsilon_0\|^2_{H^1} +  \int_{y>0} \, y^k \, \varepsilon^2_0(y)\, dy < s_0^{-1}\, , \quad (\varepsilon_0,y\Lambda Q) = (\varepsilon_0, \Lambda Q) = (\varepsilon_0, Q) = 0\, .
\end{equation} 
    Then, there exists 
    \begin{equation}\label{init_cond_in_BS}
        \lambda_0 \in [\lambda^-_0,\lambda^+_0], \quad \text{ where } \quad \lambda^{\pm}_0 := \Big( \frac{1}{\int Q}\frac{2 c_0}{1-\theta} \sigma^{1-\theta}_0
        \pm s^{-\frac{\beta}{2}-2\rho}_0 \Big)^2
        %:= \{\lambda_0\in \RR , \lvert \lambda_0 - s_0^{-\beta} \rvert \leq s_0^{-\beta-2\rho}   \}
    \end{equation}
such that the solution $U$ of \eqref{gKdV_principal_eq} evolving from the initial data
\begin{equation}
    U_0(x):= \lambda_0^{-\frac{1}{2}}\Big( Q_{b_0}+\lambda_0^{-\frac{1}{2}}f(\tau(s_0),\sigma_0)R + \vare_0 \Big)\Big( \frac{x-\sigma_0}{\lambda_0}\Big) + f(\tau(s_0),x)
\end{equation}
has a decomposition as in Lemma \ref{lemma_on_decomposition_around_Q} and satisfies \eqref{BS1}, \eqref{BS2} and \eqref{BS3} on the interval $[s_0,+\infty)$.
\end{proposition}

\begin{proof}[Proof of the Proposition \ref{Prop_on_blow_up}] 
By the continuity argument, $U$ has a decomposition as in Lemma \ref{lemma_on_decomposition_around_Q} on some small time interval $[s_0,s_1]$ with $s_1>s_0$.

For $\lambda_0 \in \mathcal{D}_0=[\lambda_0^-,\lambda_0^+]$, we define
\begin{equation}
    s^*(\lambda_0) :=\sup\{s\geq s_0, \text{ Lemma \ref{lemma_on_decomposition_around_Q}  applies and \eqref{BS1}-\eqref{BS2}-\eqref{BS3} hold on $[s_0,s]$} \}.
\end{equation}
For the sake of contradiction, we suppose for all $\lambda_0 \in \mathcal{D}_0$,
it holds $s^*(\lambda_0)< +\infty$.
The supremum $s^*$ is given only by the saturation of the bootstrap assumptions.

%By the study made in previous chapters, we know that on $[s_0,s^*]$ the estimations \todo{ref this estimations}, which we will use frequently in the following, hold true.

\vspace{0.4cm}
First, we strictly improve the bootstrap estimates \eqref{BS1} and \eqref{BS2} on $[s_0,s^*]$. Then, we find a contradiction using a continuity argument on the bootstrap assumption \eqref{BS3}.

%----------------------Closing BS3--------------------------------------------------------
\vspace{0.2cm}
\textit{Closing \eqref{BS2}.}\\
We apply \eqref{Lyapounov_control_of_H} with $j= \frac{5}{2}$ 
%(it is a technical choice to have $2<j<3$, in general I believe we suppose BS on the norm with $s^{-j}$ et on reussi a avoir une majoration avec $\frac{1}{2}s^{-j}$ et pour simplifier certains majorations et pas seulement, je crois c'est meme une contrainte, on aura besoin que $2<j<3$)
\begin{equation}
    \frac{d}{ds}\big[ \mathcal{H} \big] \;\lesssim\;  \lambda^k s^{-\frac 52} + s^{-\frac{3}{2}}.
\end{equation}
This choice $j=\frac 52$ is justified on the one hand by the need of having an error term on the 
right-hand side which is integrable in time and on the other hand by the fact that
a larger $j$ gives a better final estimate on $\vare$.
Integrating on $[s_0,s]$ for some $s\leq s^*$, we get
\begin{equation}
    s^{\frac{5}{2}}\mathcal{F}(s) \;\lesssim\; s^{\frac{5}{2}}_0 \mathcal{F}(s_0) + \lambda^k(s_0) \int \vphi_k\,\vare^2_0 + s_0^{-\frac 12}.
\end{equation}
By our choice of initial conditions in \eqref{conditons_ortho_and_decroissance_on_vare_0}, we have 
\begin{equation}
    s_0^\frac52\big| \mathcal{F}(s_0)\big| + \int \vphi_k\,\vare^2_0  \lesssim s_0^{-1}.
\end{equation}
The coercivity property \eqref{coercivity_of_mathcal_F}, \eqref{BS2} and the lower bound for $k$, yield, for $s_0$ large enough,
\begin{equation}
    \mathcal{N}^2_B(s) \;\lesssim\; \mathcal{F}(s)  \;\lesssim\; s^{-\frac 52}s_0^{-\frac12}
    \leq \frac12 s^{-\frac 52}.
\end{equation}
Thus, the estimate of $\mathcal{N}_B$ in \eqref{BS2} is strictly improved.
The control of the $H^1$ norm provided by \eqref{estimate_on_L2_norm_of_vare} and \eqref{estimate_on_L2_norm_of_grad_vare} completes the improvement of the
bootstrap estimate \eqref{BS2} for $x_0$ and $s_0$ sufficiently large.


%----------------------Closing BS2--------------------------------------------------------
\vspace{0.2cm}
\textit{Closing \eqref{BS1}.} 
\begin{itemize}
%--------BS2'_sigma-----------------------------------------------------------------------------------------------
    \item $(BS1_{\sigma})$: Using \eqref{BS1},\eqref{BS2} and the relation $(1-\beta)(1-\theta) = -\frac{\beta}{2}$, we get
\begin{equation}\label{lambda_sigma_2_2theta}
    \begin{split}
        \Big\lvert \lambda(s) - \Big( \frac{2}{\int Q}c_0 \frac{1}{\theta -1} \Big)^2 \sigma^{2-2\theta}(s) \Big\rvert \lesssim |h(s)| \Big\lvert \lambda^{\frac{1}{2}}(s) - \frac{2}{\int Q} c_0 \frac{1}{\theta -1} \sigma^{1-\theta}(s)\Big\rvert 
        \lesssim s^{-\beta-2\rho}
    \end{split}
\end{equation}
    and 
\begin{equation}
    \begin{split}
        \Big\lvert \sigma_s - \Big( \frac{2}{\int Q}c_0 \frac{1}{\theta -1} \Big)^2 \sigma^{2-2\theta}(s) \Big\rvert \lesssim |h(s)| \Big| \sigma_s^{\frac{1}{2}} - \lambda^{\frac{1}{2}} \Big| \Big\lvert \sigma_s^{\frac{1}{2}}(s) - \frac{2}{\int Q} c_0 \frac{1}{\theta -1} \sigma^{1-\theta}(s)\Big\rvert 
        \lesssim s^{-\beta-2\rho}.
\end{split}
\end{equation}
Using 
\begin{equation}\label{relation_on_constant_2theta_minus_one}
    \Big( \frac{2}{\int Q} c_0 \frac{1}{\theta -1}\Big)^2 = (1-\beta )^{-\frac{\beta}{1-\beta}} = (2\theta -1)^{2\theta - 2},
\end{equation}

yields
\begin{equation}
    \Big\lvert \big(\sigma^{\frac{1}{1-\beta}}(s)\big)_s - (1-\beta)^{-\frac{1}{1-\beta}} \Big\rvert = \frac{1}{1-\beta} \sigma^{\frac{\beta}{1-\beta}}(s) \Big\lvert \sigma_s - \Big( \frac{2}{\int Q}c_0 \frac{1}{\theta -1} \Big)^2 \sigma^{2-2\theta}(s) \Big\rvert \lesssim s^{-2 \rho}.
\end{equation}


The initial condition \eqref{init_cond_sigma_b} on $\sigma_0$, the mean value theorem for the function $x \mapsto x^{\frac{1}{1-\beta}}$, $\frac{1}{1-\beta} > 2$ and the integration of the previous result yields
\begin{equation}
    \Big\lvert \sigma^{\frac{1}{1-\beta}}(s) - (1-\beta)^{-\frac{1}{1-\beta}} s\Big\rvert \lesssim s^{1-2\rho} \iff \Bigg\lvert \Bigg( \frac{(1-\beta)\sigma(s)}{s^{1-\beta}}\Bigg)^{\frac{1}{1-\beta}} -1  \Bigg\rvert \lesssim s^{-2\rho}.
\end{equation}
By the mean value theorem as previously, we find 
\begin{equation}
    \Bigg\lvert  \frac{(1-\beta)\sigma(s)}{s^{1-\beta}}  -1  \Bigg\rvert \lesssim s^{-2\rho} \iff \Big\lvert \sigma(s) - \frac{1}{1-\beta}s^{1-\beta} \Big\rvert \lesssim s^{1-\beta-2\rho}.
\end{equation}

%----------BS2'_lambda----------------------------------------------------------------------------------------------------
\item $(BS1_\lambda)$: From \eqref{lambda_sigma_2_2theta} we get
\begin{equation}
    \Big\lvert \lambda(s) - s^{-\beta} \Big\rvert \lesssim s^{-\beta-2\rho}  + \Bigg\lvert  \Bigg( \frac{2}{\int Q}c_0 \frac{1}{\theta -1} \Bigg)^2 \sigma^{2-2\theta}(s) - s^{-\beta} \Bigg\rvert.
\end{equation}
By $(BS1_{\sigma})$, \eqref{relation_on_constant_2theta_minus_one}, the mean value theorem and since $\beta > \frac{1}{2}$, we get the desired estimate
\begin{equation}
    \Big\lvert \lambda(s) - s^{-\beta} \Big\rvert \lesssim s^{-\beta-2\rho} + s^{1-3\beta-2\rho} \lesssim s^{-\beta-2\rho}.
\end{equation}

%----------BS2'_b------------------------------------------------------------------------------------------------------
\item $(BS1_b)$: We write
\begin{equation}
    \begin{split}
        &\Big\lvert \frac{1}{\beta}sb(s) -1 \Big\rvert = \frac{1}{\beta}\lambda^2(s)s \Big\lvert \frac{b(s)}{\lambda^2(s)} - \beta s^{-1}\lambda^{-2}(s)\Big\rvert \\
        &\lesssim \lambda^2(s)s \Bigg( \lvert g(s) \rvert + \Big\lvert \frac{4}{\int Q}c_0 \lambda^{-\frac{3}{2}}(s)\sigma^{-\theta}(s) + \beta s^{-1}\lambda^{-2}(s) \Big\rvert \Bigg) .
    \end{split}
\end{equation}
 By definition of $c_0$ and of $\beta$ we find
 \begin{equation}
 \begin{split}
     &\Big\lvert \frac{4}{\int Q}c_0 \lambda^{-\frac{3}{2}}(s)\sigma^{-\theta}(s) + \beta s^{-1}\lambda^{-2}(s) \Big\rvert  = 2(\theta-1)(2\theta-1)^{-1}\lvert \lambda^{-\frac{3}{2}}(s)\rvert \Bigg\lvert \Big( \frac{\sigma(s)}{2\theta-1} \Big)^{-\theta} - s^{-1}\lambda^{-\frac{1}{2}} \Bigg\rvert\\
     &\lesssim  \lvert \lambda^{-\frac{3}{2}}(s)\rvert \Bigg( \Big\lvert \Big( \frac{\sigma(s)}{2\theta-1} \Big)^{-\theta} - s^{-1+\frac{\beta}{2}} \Big\rvert +  \lvert s^{-1}\lambda^{-\frac{1}{2}}(s)- s^{-1+\frac{\beta}{2}} \rvert \Bigg)\\
     &\lesssim s^{\frac{3}{2}\beta}(s^{-1-2\rho}+ s^{-1-2\beta-2\rho}) \lesssim s^{\frac{3}{2}\beta-1-2\rho}.
\end{split}
 \end{equation}
We have used the next result. Since $2\theta-1 = (1-\beta)^{-1}$, using $(BS1_{\sigma})$ we have
 \begin{equation}
 \begin{split}
     \Bigg\lvert \Big( \frac{\sigma(s)}{2\theta-1} \Big)^{-\theta} - s^{-1+\frac{\beta}{2}} \Bigg\rvert = \Big\lvert \Big( (1-\beta)\sigma(s) \Big)^{-\theta} - \Big(s^{1-\beta}\Big)^{-\theta}\Big\rvert \lesssim s^{-1-2\rho}.
    \end{split}
 \end{equation}
 
Thus by $(BS1_\lambda)$ we find
\begin{equation}
    \Big\lvert \frac{4}{\int Q}c_0 \lambda^{-\frac{3}{2}}(s)\sigma^{-\theta}(s) + \beta s^{-1}\lambda^{-2}(s) \Big\rvert \lesssim s^{\frac{3}{2}\beta-1-2\rho}.
\end{equation}

We conclude then, using the estimates \eqref{conseq_of_BS_on_g_and_hs} and \eqref{BS3}
\begin{equation}
\begin{split}
    &\Big\lvert \frac{1}{\beta}sb(s)-1 \Big\rvert \lesssim \lambda^2(s)\,s\,\Big(|g(s)| + \Big|\frac{4}{\int Q} c_0\,\lambda^{\frac 32}(s)\,\sigma^{-\theta}(s) + \beta\,s^{-1}\,\lambda^{-2}(s) \Big| \Big)\\ 
    &\lesssim s^{1-2\beta}\Big( s^{2\beta-1-2\rho} + s^{\frac{3}{2}\beta-1-2\rho} \Big) \lesssim s^{-2\rho} .
\end{split}
\end{equation}
\end{itemize}

%---------------------------------Topological argument--------------------------------------------
\vspace{0.2cm}
\textit{Contradiction by a continuity argument.}\\
Define the map
\begin{equation}
    \Phi: \lambda_0 \in [\lambda^-_0,\lambda^+_0] \xmapsto{\Phi_1} \mu_0 \in [-1,1] \xmapsto{\Phi_2} s^*\in [s_0,+\infty)\xmapsto{\Phi_3} h(s^*)(s^*)^{\frac{\beta}{2}+2\rho}\in \{-1,1\}.
\end{equation}
In the following, we prove that the map $\Phi$ is continuous 
and $\Phi(\lambda_0^-)=-1$, $\Phi(\lambda_0^+)=1$, which
leads to a contradiction. 

The map $\Phi_1$ is continuous and such that $\Phi_1(\lambda^{\pm}_0) = \pm1$. By the definition of $h$, the map $\Phi_3$ is also continuous. 

%We remark that $\mu_0 = \Phi(\lambda_0) = s_0^{2\rho}(s^{\beta}_0 \lambda_0 -1 ) \in [-1,1]$ and the map $\Phi_1$ is continuous. 

We define $G(s) = h^2(s)s^{\beta+4\rho}$. For $s_1 \in [s_0,s^*]$ such that $G(s_1) = 1$, the following transversality property holds true
\begin{equation}
\begin{split}
    G'(s_1) 
    &= 2h_s(s_1)h(s_1)(s_1)^{\beta + 4\rho} + (\beta + 4\rho)h^2(s_1)(s_1)^{\beta+4\rho-1}\\
    &= \pm 2 h_s(s_1)(s_1)^{\frac{\beta}{2}+2\rho} + (\beta + 4\rho)\frac{G(s_1)}{s_1}\\
    &\geq -2(s_1)^{-1-2\rho}+\frac{1}{2s_1} \geq -\frac{1}{20s_1}+ \frac{1}{2s_1} >\frac{1}{20s_1} >0.
\end{split}
\end{equation}

In particular, $G(s^*) = 1$ and $G(s_0) = 1$. From the definition of $\lambda^{\pm}_0$, the initial conditions \eqref{init_cond_sigma_b}, \eqref{conditons_ortho_and_decroissance_on_vare_0} and the transversality property on $s_0$, it holds $\Phi_2 \circ \Phi_1 (\lambda^{\pm}_0) = s_0$. Thus, $\Phi(\lambda^{\pm}_0)= \pm1$.


It remains to prove that the map $\Phi_2$ is continuous. Let $\mu_0\in (-1,1)$, so we have $s^* > s_0$. 

By continuity of $G$ and $G'$ we know that for some $\epsilon <s^*-s_0$ and $\delta_\epsilon>0$ sufficiently small holds $G(s)<1-\delta$ for all $s \in [s_0,s^*-\epsilon]$ and $G(s^*+\epsilon)> 1+\delta$. 
By continuous dependence on initial data for parameters and continuity of $\Phi_1$, there exists $\epsilon_\mu$, such that 
\begin{equation}
    \lvert \mu_0 - \tilde{\mu}_0\rvert\leq \epsilon_\mu \Rightarrow \lvert G(s) - \tilde{G}(s) \rvert \leq \frac{\delta}{2} \qquad 
    \text{on }\;[s_0,s^*+\epsilon].
\end{equation}

Let $\tilde{G}$ and $\tilde{s}^*$ be the quantities associated with $\tilde{\mu}_0$. 
We get 
\begin{equation*}
    \forall s \in [s_0,s^*-\epsilon], \quad  \tilde{G}(s) \leq G(s)+\frac{\delta}{2} < 1-\frac{\delta}{2} \quad\Rightarrow\quad \tilde{s}^* \geq s^*-\epsilon
\end{equation*}
and
\begin{equation*}
    \tilde{G}(s^*+\epsilon)\geq G(s^*+\epsilon)-\frac{\delta}{2} >1+\frac{\delta}{2} \quad\Rightarrow\quad \tilde{s}^* \leq s^*+\epsilon.
\end{equation*}
This concludes the continuity of $\Phi_2$ for $\mu_0 \in (-1,1)$. In the case $\mu_0 = \pm1$, we get $s^* = s_0$. By a similar argument, we get the continuity for $\mu_0 = \pm1$.
\end{proof}

\end{document}
